% GENERATED_BY: oc_core_monograph_translation_draft_pipeline_v1.py
% AI_ASSISTED_TRANSLATION_DRAFT: TRUE
% THEOREM_REVIEW_STATUS: PENDING
% THEOREM_REVIEW_MODE: FORMAL_AI_GATE__CODEX_CLI_CHATGPT
% THEOREM_REVIEW_REVIEWER_ID:
% THEOREM_REVIEW_AUTHORITY_CLASS:
% THEOREM_REVIEWED_AT:
% THEOREM_REVIEW_DOSSIER_REF:
% SOURCE_LANGUAGE: EN
% TARGET_LANGUAGE: RU
% SOURCE_EN_REF: content/k_levels/k2.tex
% ================================================================
% ==== FILE: content/k_levels/k2.tex
% ================================================================

% ==============================
%  Ontology of Continua — Core
%  K-level Module: K2
%  File: content/k_levels/k2.tex
%  Status: FULLY DEFINED
% ==============================



\paragraph{Canonical tuple projection note.}
Local K-level displays in this module are projections of the canonical public tuple \(K=(\Omega,\partial\Omega,A,\Theta,P,J,C,k,M)\). When a display omits \(M\), reorders components, or uses level-indexed shorthand, the omitted embedding context is inherited from the surrounding level discussion rather than introduced as a competing definition.

\subsubsection{\texorpdfstring{$K_2$}{K_2} Обзор}
\label{sec:k2-overview}

Уровень $K_2$ — это первый континуум, поддерживающий
\emph{настоящую двумерную структуру}. Он возникает через
размерностный переход
$\Psi_{1\to2}$, когда ограничения $K_1$ становятся недостаточными для
описания несовместимых вариаций (аксиома~15). Новая ось $A_2$ ортогональна
в структурном смысле и не может быть выражена как функция от $A_1$.

В результате $K_2$ становится первым уровнем, на котором:
\begin{itemize}
    \item существуют нетривиальные циклы,
    \item возникает время $\tau$ (через устойчивые $C$-циклы),
    \item градиенты и потенциалы формируют двумерные поля,
    \item появляются дефекты, вихри и дислокации,
    \item топологические явления приобретают смысл.
\end{itemize}

$K_2$ — минимальная среда для классической поведенческой
структуры полевых теорий и перехода типа BKT.

% ================================================================
\subsubsection{\texorpdfstring{$\Omega(K_2)$}{\Omega(K_2)} Пространство состояний}

Пространство состояний $K_2$:
\[
\Omega(K_2) = C^0(X_1 \times X_2, V),
\]
пространство непрерывных полей на 2D-области.

Ключевые свойства:
\begin{enumerate}
    \item $\Omega(K_2)$ содержит конфигурации со значимыми локальными
          градиентами по двум независимым направлениям.
    \item Дефекты (например, вихри) лежат в $\partial\Omega(K_2)$,
          если они не удовлетворяют ограничениям допустимости из $M_2$.
    \item Существуют нетривиальные гомотопические классы:

          \[
          \pi_1(\Omega(K_2)) \neq 0.
          \]
\end{enumerate}

Это первый уровень, где возникает геометрия (в слабом структурном смысле).

% ================================================================
\subsubsection{\texorpdfstring{$\partial\Omega(K_2)$}{\partial\Omega(K_2)} Граница}

$\partial\Omega(K_2)$ содержит:
\begin{itemize}
    \item разрывные поля,
    \item поля с расходящейся плотностью энергии,
    \item конфигурации, в которых превышаются $\Theta_{\mathrm{grad}}$ или
          $\Theta_{\mathrm{defect}}$,
    \item поля, где циклы коллапсируют ($C \to \varnothing$).
\end{itemize}

Специальная часть границы соответствует топологическим дефектам,
ядро которых схлопывается до нуля, что указывает на разрушение
размерностной согласованности.

% ================================================================
\subsubsection{\texorpdfstring{$A(K_2)$}{A(K_2)} Оси}

Оси состоят из:
\[
A(K_2) = \{A_1, A_2\},
\]
где:
\begin{itemize}
    \item $A_1$ наследуется от $K_1$,
    \item $A_2$ — новая независимая структурная ось.
\end{itemize}

Условие независимости:
\[
\nexists\, h: \mathbb{R} \to \mathbb{R} \\text{ при этом } A_2 = h \circ A_1.
\]

Эта независимость является формальным ядром размерностного скачка.

% ================================================================
\subsubsection{\texorpdfstring{$P(K_2)$}{P(K_2)} Потенциалы}

Множество потенциалов теперь включает полные 2D-структурные члены:
\[
P(K_2) =
\{ P_{\mathrm{grad}}, P_{\mathrm{curl}}, P_{\mathrm{div}},
   P_{\mathrm{bond}}, P_{\mathrm{smooth}}, P_{\mathrm{top}} \}.
\]

Интерпретация:
\begin{itemize}
    \item градиенты по двум осям кодируют локальное напряжение,
    \item вихревые структуры кодируют вращательные степени свободы,
    \item потенциалы дивергенции кодируют сжимающее или расширительное
          напряжение,
    \item $P_{\mathrm{top}}$ описывает энергию вихрей и дефектов.
\end{itemize}

Эти потенциалы воспроизводят на абстрактном уровне структурные
предшественники полей в физике.

% ================================================================
\subsubsection{\texorpdfstring{$\Theta(K_2)$}{\Theta(K_2)} Пороги}

Пороги теперь включают:
\[
\Theta(K_2) =
\{\Theta_{\mathrm{grad}}, \Theta_{\mathrm{curl}},
  \Theta_{\mathrm{defect}}, \Theta_{\mathrm{coh}},
  \Theta_{\mathrm{BKT}}\}.
\]

Их смысл:
\begin{itemize}
    \item $\Theta_{\mathrm{grad}}$: максимальная допустимая величина градиента.
    \item $\Theta_{\mathrm{curl}}$: допустимая кривизна вращения.
    \item $\Theta_{\mathrm{defect}}$: порог устойчивости дефектного ядра.
    \item $\Theta_{\mathrm{coh}}$: порог согласованности для 2D-упорядоченности.
    \item $\Theta_{\mathrm{BKT}}$: порог появления устойчивых $C$-циклов;
          эквивалентен $K_c$ в теореме представимости BKT.
\end{itemize}

Пересечение $\Theta_{\mathrm{BKT}}$ даёт ось возникающего времени.

% ================================================================
\subsubsection{\texorpdfstring{$J(K_2)$}{J(K_2)} Потоки}

Потоки включают:
\[
J(K_2) = \{\partial_{x_1} f,\; \partial_{x_2} f,\;
\text{2D деформации},\; \text{вихревое дрейфование},\; \text{повороты фазы}\}.
\]

Здесь появляются:
\begin{itemize}
    \item дивергентные и вихревые потоки,
    \item диффузионное выравнивание,
    \item динамика аннигиляции дефектов и разъединения дефектных пар.
\end{itemize}

Потоки задают динамическую допустимость эволюции в $\Omega(K_2)$.

% ================================================================
\subsubsection{\texorpdfstring{$C(K_2)$}{C(K_2)} Циклы}

$K_2$ — первый уровень с нетривиальными циклами:
\[
C(K_2) = \{C_{\mathrm{vortex}},\; C_{\mathrm{phase}},\; C_{\mathrm{ring}}\}.
\]

Самый важный:
\[
C_{\mathrm{vortex}}: \oint \nabla \theta \cdot d\ell = 2\pi n.
\]

Согласно Теореме Представимости~5 (BKT):
\begin{itemize}
    \item Для связи $K < K_c$ циклы $C_{\mathrm{vortex}}$ неустойчивы и
          распадаются.
    \item Для $K > K_c$ циклы становятся устойчивыми и обеспечивают
          глобальную согласованность.
\end{itemize}

Эти устойчивые циклы порождают временную степень свободы.

% ================================================================
\subsubsection{\texorpdfstring{$\tau(K_2)$}{\tau(K_2)} Время}

Время возникает в $K_2$ впервые.

Из теоремы происхождения времени:
\[
\tau \text{ существует тогда и только тогда, когда есть устойчивый цикл }
C \text{ с } \Pi(C) > \Theta_{\mathrm{time}}.
\]

Здесь $\Pi(C)$ — функционал мощности цикла.

Следовательно:
\[
\tau(K_2) =
\begin{cases}
\text{не определено}, & K < K_c, \\
\text{хорошо определённый монотонный параметр}, & K > K_c.
\end{cases}
\]

Это структурный аналог возникновения временной упорядоченности из
устойчивой двумерной согласованности.

% ================================================================
\subsubsection{\texorpdfstring{$k(K_2)$}{k(K_2)} Континуумность}

Общая формула:
\[
k(K_2) =
\frac{\mu(\Omega(K_2))}{\mu(S(K_2))}
\cdot
\frac{|A(K_2)|}{|A(M_2)|}
\cdot
\frac{\sum \mathrm{Stab}(C)}{\sum \mathrm{MaxStab}(C)}
\cdot
\left(1 - \frac{\sigma(J)}{\sigma_{\max}}\right)
\cdot
\left(1 - \frac{T}{\Theta_{\mathrm{coh}}}\right)_+.
\]

Интерпретация:
\begin{itemize}
    \item рост $\Omega(K_2)$ увеличивает $k$,
    \item повышение плотности дефектов уменьшает $k$,
    \item устойчивые циклы повышают $k$,
    \item турбулентные потоки уменьшают $k$.
\end{itemize}

Это количественно описывает структурную жизнеспособность двумерного континуума.

% ================================================================
\subsubsection{\texorpdfstring{$T(K_2)$}{T(K_2)} Структурное напряжение}

Источники напряжения:
\begin{itemize}
    \item крупные градиенты,
    \item ядра вихрей,
    \item взаимодействия дефектов-антиизоморфов,
    \item несовместимые краевые условия.
\end{itemize}

Коллапс наступает при:
\[
T(K_2) > \Theta_{\mathrm{coh}}.
\]

Это приводит к потере согласованности и разрушению $C$-циклов.

% ================================================================
\subsubsection{\texorpdfstring{$E(K_2)$}{E(K_2)} Энергия}

Структурная энергия:
\[
E(K_2) = \int_{X_1 \times X_2}
\mathcal{E}(f, \nabla f, \nabla^2 f)\, d\mu,
\]
где $\mathcal{E}$ включает:
\begin{itemize}
    \item градиентную энергию,
    \item эластоподобную энергию,
    \item энергию сердцевин дефектов,
    \item топологическую энергию вихрей.
\end{itemize}

Это структурный прототип энергетических функционалов в физике.

% ================================================================
\subsubsection{\texorpdfstring{$K_2$}{K_2} Операторы (\texorpdfstring{$\Psi$}{\Psi}, \texorpdfstring{$\Phi$}{\Phi}, \texorpdfstring{$\Lambda$}{\Lambda}, \texorpdfstring{$U$}{U}, \texorpdfstring{$\Chi$}{\Chi})}

\begin{itemize}
    \item $\Psi_{2\to3}$ — размерностный переход к $K_3$, когда
          двумерная структура становится недостаточной.
    \item $\Phi$ — оператор эволюции мета-пространства для $M_2$.
    \item $\Lambda$ — оператор жёсткого ограничения, стабилизирующий
          градиенты, гладкость и границы дефектов.
    \item $U$ — унификация локальных 2D-областей в глобальные согласованные
          поля.
    \item $\Chi$ — ветвящийся оператор, допускающий несколько многообразий
          $K_2$ с разной структурой согласованности.
\end{itemize}

% ================================================================
\subsubsection{Процессы на \texorpdfstring{$K_2$}{K_2}}

Процессы включают:
\begin{itemize}
    \item создание, аннигиляцию и дрейф вихрей,
    \item фазовую упорядоченность и дезорганизацию,
    \item переходы типа BKT,
    \item стабилизацию глобальной согласованности,
    \item коллапс при росте дефектов,
    \item размерностной рост к $K_3$, когда несовместимые вариации
          нельзя описать через $\{A_1, A_2\}$.
\end{itemize}

% ================================================================
\subsubsection{Прогнозы для \texorpdfstring{$K_2$}{K_2}}

\begin{enumerate}
    \item Наличие порога согласованности $K_c$ (BKT).
    \item Резкий переход между устойчивыми и неустойчивыми вихрями.
    \item Возникновение времени только при стабилизации вихревых циклов.
    \item Универсальные степенные корреляции ниже $K_c$.
    \item Потеря согласованности при росте дефектной плотности.
    \item Необходимость двух независимых осей для нетривиальных циклов.
\end{enumerate}

Эти предсказания согласуются с широким классом физических систем (модель XY,
сверхтекучие плёнки, тонкие магнитные слои).

% ================================================================
\subsubsection{Эксперименты для \texorpdfstring{$K_2$}{K_2}}

Канонические экспериментальные сигнатуры:
\begin{itemize}
    \item измерение температуры расцепления вихрей в двумерных сверхтекучих средах,
    \item обнаружение BKT-подобных переходов в атомарно тонких пленках,
    \item наблюдение коллапса согласованности при высокой плотности дефектов,
    \item кинетика фазового упорядочения в 2D-материалах,
    \item универсальный скачок жесткости на BKT-переходе.
\end{itemize}

Косвенные проверки:
\begin{itemize}
    \item динамика дефектов в графене,
    \item 2D-пленки жидких кристаллов,
    \item биологические мембраны с топологическими паттернами дефектов.
\end{itemize}

% ================================================================
\subsubsection{Смерть и коллапс \texorpdfstring{$K_2$}{K_2}}

$K_2$ коллапсирует, когда:
\[
\Omega(K_2) = \varnothing,
\]
по причине:
\begin{itemize}
    \item неконтролируемого роста дефектов,
    \item потери согласованности ($T > \Theta_{\mathrm{coh}}$),
    \item разрушения всех $C$-циклов,
    \item невозможности поддерживать две независимые оси.
\end{itemize}

После коллапса переход к $K_3$ невозможен.

% ================================================================
\subsubsection{Фальсифицируемость \texorpdfstring{$K_2$}{K_2}}

Фальсифицируемые предсказания:
\begin{enumerate}
    \item существование резкого BKT-подобного порога,
    \item необходимость циклов для возникающего времени,
    \item поведение корреляционной длины около $K_c$,
    \item универсальный скачок жесткости,
    \item топологическая защита $C_{\mathrm{vortex}}$ выше $K_c$.
\end{enumerate}

Эмпирическое нарушение этих признаков опровергнет структурную
модель.

% ================================================================
\subsubsection{Ветвление / Онтологическая позиция \texorpdfstring{$K_2$}{K_2}}

$K_2$ является корнем всех континуумов более высоких размерностей:
\[
K_0 \to K_1 \to K_2 \to K_3 \to \dots
\]

Ветвление на $K_2$ соответствует различным допустимым двумерным
структурам согласованности:
\begin{itemize}
    \item фазы с большим количеством вихрей и без вихрей,
    \item гладкие и дефектно-ориентированные многообразия,
    \item ветви с высокой и низкой согласованностью.
\end{itemize}

Эти ветви остаются различными, пока их не объединяет оператор $U$.

% ================================================================
\subsubsection{Связь с M-пространствами}

$K_2$ вложен:
\[
K_2 \subseteq M_2 \subseteq M_1 \subseteq M_0.
\]

$M_2$ определяет:
\begin{itemize}
    \item допустимые 2D топологии,
    \item допускаемые классы дефектов,
    \item ограничения градиента и роторного поля,
    \item условия согласованности для размерностной устойчивости.
\end{itemize}

Переход к $K_3$ требует, чтобы $M_2$ не мог вместить новые
несовместимые классы различий.
